\section{Resultados}
\label{sec:resultados}

Los resultados se presentan en el mismo orden que los protocolos. Primero se
resumen las reconstrucciones bibliográficas; después se muestran los dos casos
obtenidos por la metodología propuesta. En el caso no suave se completa el
primer radio con contactos. Para \texttt{c590} se distingue la partición local
\(r\leq0.3\) de la auditoría macroscópica en \(r=1\) y \(r=2\).

\subsection{Reconstrucciones bibliográficas}
\label{subsec:resultados_bibliograficos}

\subsubsection{Recuperación del caso entero}

Para \cref{eq:parametros_enteros}, la primera rama del cierre de
Leonov--Kuznetsov produjo
\begin{equation}
\omega_0=2.039186939959,\qquad
k=0.209867354515,\qquad
A_0=5.856145086257,
\label{eq:resultados_armonicos_enteros}
\end{equation}
y la condición modal
\begin{equation}
X_0=(5.856145086257,\ 0.369331578247,\ -8.366536168332)^{\T}.
\label{eq:resultado_semilla_entera}
\end{equation}
Los valores propios de la matriz auxiliar fueron
\[
\lambda_1=-1.538510163250,\qquad
\lambda_{2,3}\simeq\pm2.039186939959\,\mathrm i,
\]
con residuo complejo de cierre \(9.93\times10^{-16}\). Bajo el carril
histórico EFORK--3 \((T=320,\ T_{\mathrm{burn}}=80)\) y su clasificador por
huellas de sección, las 504 sondas produjeron 369 convergencias a equilibrio,
69 divergencias y 66 trayectorias clasificadas como otro destino; ninguna hizo
contacto con la nube objetivo.

\subsubsection{Danca y Wu bajo sus contratos reconstruidos}

La reconstrucción nominal de Danca con
\cref{eq:semilla_danca_diagnostica} permaneció acotada. Después del transitorio
presentó una frecuencia dominante de \(0.3679852806\) ciclos por unidad de
tiempo y medianas de la prueba 0--1
\begin{equation}
(K_x,K_y,K_z)
=(0.0736852160,\ 0.0214448125,\ -0.0012988924).
\label{eq:resultado_danca_k}
\end{equation}
Estos valores corresponden a una trayectoria de apariencia periódica bajo el
contrato ABM de \cref{subsec:protocolos_bibliograficos}.

La recurrencia ADM de Wu permaneció acotada para las dos condiciones de
\cref{eq:semillas_wu_protocolo}. Las trayectorias fueron simétricas y de
apariencia periódica después del transitorio. Las tres componentes tuvieron
frecuencia dominante \(0.5398920216\) ciclos por unidad de tiempo, deriva
relativa nula y entropías espectrales
\[
(H_x,H_y,H_z)
=(0.11607348,\ 0.11106582,\ 0.11095067).
\]

\begin{table}[H]
\centering
\caption{Resultados numéricos de los tres controles bibliográficos.}
\label{tab:resultados_referencias}
\begin{tabularx}{\linewidth}{L{0.20\linewidth}L{0.27\linewidth}Y}
\toprule
Control & Configuración evaluada & Resultado reproducido\\
\midrule
Leonov--Kuznetsov, carril histórico EFORK--3
& \(q=1\), parámetros de \cref{eq:parametros_enteros}
& \(\omega_0=2.03918694\), \(A_0=5.85614509\);
0 contactos en 504 sondas\\
Danca
& \(q=0.9998\), ABM completo, \(h=0.01\), \(T=500\)
& Trayectoria acotada de apariencia periódica,
\(f_d=0.36798528\)\\
Wu et al.
& \(q=0.99\), ADM local de orden cuatro, dos semillas publicadas
& Dos trayectorias simétricas de apariencia periódica,
\(f_d=0.53989202\)\\
\bottomrule
\end{tabularx}
\end{table}

\begin{hafodigitalfigure}[H]
\centering
\includegraphics[width=0.96\linewidth]{bibliographic_reproducibility_controls_es}
\caption{Controles de reproducibilidad. Los paneles de Danca y Wu muestran las
trayectorias obtenidas con sus contratos reconstruidos. La continuación no
suave propuesta y el control con pendientes distintas se presentan como
experimentos independientes y no como resultados de las referencias.}
\label{fig:controles_bibliograficos}
\end{hafodigitalfigure}

\FloatBarrier
\subsection{Caso no suave localizado con semilla sesgada}
\label{subsec:resultados_no_suave}

\subsubsection{Raíz, semilla y continuación}

La exploración de 2\,430 inicios seleccionó, dentro del par nominal,
\begin{equation}
\begin{aligned}
A&=4.577882721023,&
c&=2.776397003274,\\
\omega&=2.040286051079,&
N_1&=0.210022792962,&
\Psi_0&=0.408770376181.
\end{aligned}
\label{eq:resultado_raiz_no_suave}
\end{equation}
La norma del residuo de la cuadratura trapezoidal usada por el solucionador fue
\(4.55\times10^{-16}\). Una reevaluación independiente de las integrales
continuas mediante cuadratura adaptativa, separando los puntos de quiebre de
la saturación, produjo \(1.9\times10^{-7}\), menor que el umbral de aceptación
\(10^{-4}\). Ambas cifras corresponden a evaluaciones numéricas diferentes del
mismo cierre.

La condición inicial construida fue
\begin{equation}
X_{\mathrm{sem}}
=(7.354279724297,\ 0.290968778784,\ -9.316054818666)^{\T}.
\label{eq:resultado_semilla_no_suave}
\end{equation}
Después de transportar una sola memoria a través de los 21 niveles, el estado
final fue
\begin{equation}
X_{\mathrm{end}}
=(0.822310709850,\ 1.465272703995,\ 1.292308250248)^{\T}.
\label{eq:resultado_endpoint_no_suave}
\end{equation}

\begin{hafodigitalfigure}[H]
\centering
\includegraphics[width=0.90\linewidth]{nonsmooth_numerical_continuation}
\caption{Continuación numérica del sistema auxiliar al sistema no suave
original. Los niveles \(\varepsilon=0,0.05,\ldots,1\) comparten una sola
historia de Caputo.}
\label{fig:continuacion_no_suave_paper07}
\end{hafodigitalfigure}

El problema autónomo iniciado en \cref{eq:resultado_endpoint_no_suave}
produjo el conjunto postransitorio archivado para el caso no suave. Su parte
postransitoria tuvo apariencia periódica, frecuencia dominante
\(0.3699815009\) y entropías
\begin{equation}
(H_x,H_y,H_z)
=(0.11526785,\ 0.10497088,\ 0.10339964).
\label{eq:entropias_no_suave}
\end{equation}
En conjunto, estos diagnósticos lo clasifican como un régimen regular de
apariencia periódica, no como un atractor caótico.
Las 36 perturbaciones alrededor de \(X_{\mathrm{end}}\) hicieron contacto con
la nube objetivo: 12 de 12 para cada radio \(10^{-6}\), \(10^{-4}\) y
\(10^{-2}\).

\begin{hafodigitalfigure}[H]
\centering
\includegraphics[width=0.86\linewidth]{nonsmooth_final_attractor}
\caption{Conjunto atrayente postransitorio obtenido para
\cref{eq:chua_no_suave,eq:parametros_no_suave} desde el extremo de la
continuación causal.}
\label{fig:atractor_no_suave_paper07}
\end{hafodigitalfigure}

\subsubsection{Sondeo volumétrico}

El protocolo completo reunió 17\,400 problemas de valor inicial, 5\,800
alrededor de cada equilibrio, sin fallos numéricos. De ellos, 7\,200
corresponden al contrato local predeclarado \(r\leq0.01\), y 10\,200 a la
extensión macroscópica diagnóstica en \(r=0.03,0.1,0.3\). La
\cref{tab:sondeo_no_suave_completo} contiene todos los radios ejecutados y
todas las clases del clasificador.

\begingroup
\setlength{\tabcolsep}{10pt}
\setlength{\LTleft}{\fill}
\setlength{\LTright}{\fill}
\begin{longtable}{@{}ccrrrrr@{}}
\caption{Conteos completos del sondeo volumétrico del sistema no suave.}
\label{tab:sondeo_no_suave_completo}\\
\toprule
Equilibrio & \(r\) & \(N\) & Objetivo & Equilibrio & Divergencia & Otro\\
\midrule
\endfirsthead
\multicolumn{7}{c}{\tablename\ \thetable\ (continuación)}\\
\toprule
Equilibrio & \(r\) & \(N\) & Objetivo & Equilibrio & Divergencia & Otro\\
\midrule
\endhead
\midrule
\multicolumn{7}{r}{Continúa en la página siguiente}\\
\endfoot
\bottomrule
\endlastfoot
\multirow{10}{*}{\(E_0\)}
& \(10^{-5}\) & 80 & 0 & 80 & 0 & 0\\
& \(3{\times}10^{-5}\) & 120 & 0 & 120 & 0 & 0\\
& \(10^{-4}\) & 160 & 0 & 160 & 0 & 0\\
& \(3{\times}10^{-4}\) & 240 & 0 & 240 & 0 & 0\\
& \(10^{-3}\) & 400 & 0 & 400 & 0 & 0\\
& \(3{\times}10^{-3}\) & 600 & 0 & 600 & 0 & 0\\
& \(10^{-2}\) & 800 & 0 & 800 & 0 & 0\\
& \(0.03\) & 1000 & 0 & 1000 & 0 & 0\\
& \(0.1\) & 1200 & 0 & 1200 & 0 & 0\\
& \(0.3\) & 1200 & 0 & 1200 & 0 & 0\\
\midrule
\multirow{10}{*}{\(E_+\)}
& \(10^{-5}\) & 80 & 0 & 0 & 39 & 41\\
& \(3{\times}10^{-5}\) & 120 & 0 & 0 & 63 & 57\\
& \(10^{-4}\) & 160 & 0 & 0 & 80 & 80\\
& \(3{\times}10^{-4}\) & 240 & 0 & 0 & 116 & 124\\
& \(10^{-3}\) & 400 & 0 & 0 & 189 & 211\\
& \(3{\times}10^{-3}\) & 600 & 0 & 0 & 277 & 323\\
& \(10^{-2}\) & 800 & 0 & 0 & 393 & 407\\
& \(0.03\) & 1000 & 0 & 0 & 490 & 510\\
& \(0.1\) & 1200 & 0 & 4 & 588 & 608\\
& \(0.3\) & 1200 & 19 & 7 & 591 & 583\\
\midrule
\pagebreak[4]
\multirow{10}{*}{\(E_-\)}
& \(10^{-5}\) & 80 & 0 & 0 & 41 & 39\\
& \(3{\times}10^{-5}\) & 120 & 0 & 0 & 56 & 64\\
& \(10^{-4}\) & 160 & 0 & 0 & 81 & 79\\
& \(3{\times}10^{-4}\) & 240 & 0 & 0 & 122 & 118\\
& \(10^{-3}\) & 400 & 0 & 0 & 204 & 196\\
& \(3{\times}10^{-3}\) & 600 & 0 & 0 & 309 & 291\\
& \(10^{-2}\) & 800 & 0 & 0 & 394 & 406\\
& \(0.03\) & 1000 & 0 & 0 & 488 & 512\\
& \(0.1\) & 1200 & 0 & 4 & 593 & 603\\
& \(0.3\) & 1200 & 18 & 12 & 596 & 574\\
\end{longtable}
\endgroup

En el contrato local se registraron cero contactos y cero fallos numéricos en
las 7\,200 sondas con \(r\leq0.01\). Este es el alcance usado para evaluar la
compatibilidad local con ocultedad del conjunto atrayente regular.

La extensión macroscópica, que no amplía ese contrato local, añadió 10\,200
sondas. Las 6\,600 correspondientes a \(r=0.03\) y \(r=0.1\) produjeron cero
contactos; al combinarlas con el contrato local se obtiene el acumulado
\(0/13\,800\) hasta \(r=0.1\). El primer radio con contactos de la extensión
fue \(r=0.3\); después de completar sus 3\,600 sondas se registraron 37
contactos: cero desde \(E_0\), 19 desde \(E_+\) y 18 desde \(E_-\). Las
distancias iniciales de esos contactos a su equilibrio estuvieron entre
\(0.1941805298\) y \(0.2997254735\). Por tanto, la extensión macroscópica
delimita el intervalo de acceso muestreado
\[
0.1<r_{\mathrm{contacto}}\leq0.3
\]
sin convertirlo en el alcance del contrato local.

\begin{hafodigitalfigure}[H]
\centering
\begin{subfigure}[t]{0.49\linewidth}
  \centering
  \includegraphics[width=\linewidth]{nonsmooth_first_harmonic_comparison_es}
  \caption{Trayectoria y primer armónico.}
\end{subfigure}\hfill
\begin{subfigure}[t]{0.49\linewidth}
  \centering
  \includegraphics[width=\linewidth]{nonsmooth_volumetric_probe_overview_es}
  \caption{Sondas volumétricas representativas.}
\end{subfigure}
\caption{Diagnósticos del caso no suave. La comparación armónica describe el
carácter regular de la respuesta; las sondas muestran ejemplos de las clases
usadas en \cref{tab:sondeo_no_suave_completo}.}
\label{fig:diagnosticos_no_suave_paper07}
\end{hafodigitalfigure}

\FloatBarrier
\subsection{Caso con arctangente \texttt{c590}}
\label{subsec:resultados_c590}

\subsubsection{Selección en orden entero y refinamiento de Caputo}

En la exploración global, 894 de los 2\,400 casos quedaron separados del
control y 884 fueron además no triviales. De éstos, 823 se clasificaron como
regulares de apariencia periódica y 61 permanecieron en la categoría
exploratoria no periódica. El índice global de base cero 1731 obtuvo la mayor
puntuación, \(S=0.3021129950\), y definió el centro
\begin{equation}
\begin{aligned}
p_b={}&(18.4857295104,\ 21.9669521100,\ 0.0052746106,\\
&0.0262974930,\ -3.2087779245,\ 1.9492612035),\\
X_b={}&(7.0796053271,\ 0.5079327671,\ -14.4903935267)^{\T}.
\end{aligned}
\label{eq:resultado_centro_c590}
\end{equation}

La exploración local generó 1\,000 casos: 170 quedaron separados del control
y 148 fueron también no triviales. Entre estos últimos hubo 87 respuestas
exploratorias no periódicas y 61 regulares de apariencia periódica.
\texttt{c590}, índice local de base cero 590, ocupó el rango 10 por puntuación.
La auditoría variacional de los 30 mejores casos seleccionó \texttt{c590} con
\begin{equation}
(\lambda_1,\lambda_2,\lambda_3)
=(0.4699043683,\ 0.0010414594,\ -19.0491693716)
\label{eq:resultado_lyapunov_entero_c590}
\end{equation}
y semilla entera
\begin{equation}
X_0^{\mathrm{int}}
=(7.6733768929,\ 0.5079327671,\ -14.4903935267)^{\T}.
\label{eq:resultado_semilla_entera_c590}
\end{equation}

Los parámetros finales son los de \cref{eq:parametros_c590}. Sus equilibrios
son
\begin{equation}
E_0=(0,0,0),\qquad
E_\pm=(\pm4.6494055135,\ \pm0.0017970023,\ \mp4.6476085112).
\label{eq:equilibrios_c590_resultados}
\end{equation}
El criterio de Matignon clasificó \(E_0\) como inestable y \(E_\pm\) como
estables; el margen de los equilibrios externos fue
\(7.42265\times10^{-4}\).

El refinamiento seleccionó \(q=0.9999\) y el estado de índice 9, extraído en
\(t\simeq240\):
\begin{equation}
X_0^{\mathrm{frac}}
=(5.8642449791,\ 1.5847111486,\ 3.2155806478)^{\T}.
\label{eq:resultado_semilla_fraccionaria_c590}
\end{equation}
En la auditoría larga, las medianas remuestreadas fueron
\begin{equation}
(K_{0.0025},K_{0.005},K_{0.01})
=(0.865843,\ 0.885089,\ 0.740058).
\label{eq:resultado_k_pasos_c590}
\end{equation}
Los dos pasos refinados superaron el umbral de selección 0.8; el tercer valor
documenta la sensibilidad al paso más grande.

\begin{hafodigitalfigure}[H]
\centering
\includegraphics[width=0.92\linewidth]{arctan_fractional_refinement_es}
\caption{Búsqueda en orden entero y refinamiento fraccionario del caso
\texttt{c590}. Cada valor de \(q\) corresponde a un problema de Caputo
independiente; el estado seleccionado se extrae de la trayectoria
\(q=0.9999\).}
\label{fig:refinamiento_c590_paper07}
\end{hafodigitalfigure}

La trayectoria final permaneció acotada en los tres pasos. Sobre 30\,000
estados postransitorios, ninguna componente cumplió el filtro de apariencia
periódica. Las entropías espectrales, fracciones de potencia dominante y
derivas relativas fueron
\begin{equation}
\begin{aligned}
(H_x,H_y,H_z)
&=(0.394751,\ 0.261922,\ 0.315229),\\
(R_x,R_y,R_z)
&=(0.122163,\ 0.290492,\ 0.241937),\\
(\Delta_{f,x},\Delta_{f,y},\Delta_{f,z})
&=(0.941176,\ 0.029412,\ 0.029412).
\end{aligned}
\label{eq:metricas_finales_c590}
\end{equation}
La respuesta se clasifica, por tanto, como aperiódica bajo el filtro
postransitorio declarado.

\begin{hafodigitalfigure}[H]
\centering
\includegraphics[width=0.86\linewidth]{arctan_final_attractor}
\caption{Conjunto postransitorio del sistema
\cref{eq:chua_arctan,eq:parametros_c590} para \(q=0.9999\),
\(h=0.005\) y la condición \cref{eq:resultado_semilla_fraccionaria_c590}.}
\label{fig:atractor_c590_paper07}
\end{hafodigitalfigure}

\subsubsection{Sondeo superficial}

El registro canónico reúne 13\,500 trayectorias: 8\,400 en la partición local
\(r\leq0.3\) y 5\,100 en la auditoría macroscópica \(r\in\{1,2\}\). El primer
contacto se observó en \(r=1\), pero se retuvo \(r=2\) como bloque adicional.
La \cref{tab:sondeo_c590_completo} contiene todos los conteos del experimento
declarado.

\begingroup
\setlength{\tabcolsep}{10pt}
\setlength{\LTleft}{\fill}
\setlength{\LTright}{\fill}
\begin{longtable}{@{}ccrrrrr@{}}
\caption{Conteos completos del sondeo sobre superficies esféricas de
\texttt{c590}.}
\label{tab:sondeo_c590_completo}\\
\toprule
Equilibrio & \(r\) & \(N\) & Objetivo & Equilibrio & Divergencia & Otro\\
\midrule
\endfirsthead
\multicolumn{7}{c}{\tablename\ \thetable\ (continuación)}\\
\toprule
Equilibrio & \(r\) & \(N\) & Objetivo & Equilibrio & Divergencia & Otro\\
\midrule
\endhead
\midrule
\multicolumn{7}{r}{Continúa en la página siguiente}\\
\endfoot
\bottomrule
\endlastfoot
\multirow{9}{*}{\(E_0\)}
& \(10^{-5}\) & 100 & 0 & 0 & 0 & 100\\
& \(10^{-4}\) & 200 & 0 & 0 & 0 & 200\\
& \(10^{-3}\) & 300 & 0 & 0 & 0 & 300\\
& \(10^{-2}\) & 400 & 0 & 0 & 0 & 400\\
& \(0.03\) & 500 & 0 & 0 & 0 & 500\\
& \(0.1\) & 600 & 0 & 0 & 0 & 600\\
& \(0.3\) & 700 & 0 & 0 & 0 & 700\\
& \(1\) & 800 & 22 & 0 & 0 & 778\\
& \(2\) & 900 & 289 & 0 & 67 & 544\\
\midrule
\pagebreak[4]
\multirow{9}{*}{\(E_+\)}
& \(10^{-5}\) & 100 & 0 & 2 & 0 & 98\\
& \(10^{-4}\) & 200 & 0 & 0 & 0 & 200\\
& \(10^{-3}\) & 300 & 0 & 0 & 0 & 300\\
& \(10^{-2}\) & 400 & 0 & 0 & 0 & 400\\
& \(0.03\) & 500 & 0 & 0 & 0 & 500\\
& \(0.1\) & 600 & 0 & 0 & 0 & 600\\
& \(0.3\) & 700 & 0 & 0 & 0 & 700\\
& \(1\) & 800 & 0 & 0 & 0 & 800\\
& \(2\) & 900 & 150 & 0 & 37 & 713\\
\midrule
\multirow{9}{*}{\(E_-\)}
& \(10^{-5}\) & 100 & 0 & 2 & 0 & 98\\
& \(10^{-4}\) & 200 & 0 & 0 & 0 & 200\\
& \(10^{-3}\) & 300 & 0 & 0 & 0 & 300\\
& \(10^{-2}\) & 400 & 0 & 0 & 0 & 400\\
& \(0.03\) & 500 & 0 & 0 & 0 & 500\\
& \(0.1\) & 600 & 0 & 0 & 0 & 600\\
& \(0.3\) & 700 & 0 & 0 & 0 & 700\\
& \(1\) & 800 & 0 & 0 & 0 & 800\\
& \(2\) & 900 & 149 & 0 & 27 & 724\\
\end{longtable}
\endgroup

Las 8\,400 trayectorias con \(r\leq0.3\) produjeron cero contactos; 8\,396
completaron el horizonte y cuatro convergieron anticipadamente a \(E_\pm\) en
\(r=10^{-5}\). En \(r=1\) hubo 22 contactos entre 2\,400 trayectorias, todos
desde \(E_0\). En \(r=2\) hubo 588 contactos entre 2\,700 trayectorias: 289
desde \(E_0\), 150 desde \(E_+\) y 149 desde \(E_-\); en este bloque se
registraron 131 divergencias. En total se ejecutaron 13\,500 pruebas y se
observaron 610 contactos. El primer radio observado satisface
\[
0.3<r_{\mathrm{primer\ contacto}}\leq1.
\]

\begin{hafodigitalfigure}[H]
\centering
\begin{subfigure}[t]{0.49\linewidth}
  \centering
  \includegraphics[width=\linewidth]{arctan_first_harmonic_comparison_es}
  \caption{Trayectoria y primer armónico.}
\end{subfigure}\hfill
\begin{subfigure}[t]{0.49\linewidth}
  \centering
  \includegraphics[width=\linewidth]{arctan_spherical_probe_overview_es}
  \caption{Sondas superficiales representativas.}
\end{subfigure}
\caption{Diagnósticos de \texttt{c590}. La comparación con el primer armónico
acompaña el filtro de aperiodicidad. El panel de sondas muestra nueve historias
representativas seleccionadas en el subconjunto hasta \(r=1\); no es una
gráfica de conteos. La \cref{tab:sondeo_c590_completo} incorpora además la
auditoría canónica en \(r=2\).}
\label{fig:diagnosticos_c590_paper07}
\end{hafodigitalfigure}

\FloatBarrier
\subsection{Delimitación local obtenida}
\label{subsec:delimitacion_local}

\begin{table}[H]
\centering
\caption{Alcance local predeclarado y acceso observado en los sondeos
extendidos.}
\label{tab:intervalos_locales}
\begin{tabularx}{\linewidth}{L{0.26\linewidth}C{0.18\linewidth}C{0.18\linewidth}Y}
\toprule
Sistema & Geometría & Alcance o intervalo & Conteos correspondientes\\
\midrule
No suave, \(q=0.9998\)
& Bola volumétrica
& Local: \(r\leq0.01\)
& Local: \(0/7\,200\), cero fallos;
macro: \(0/6\,600\) en \(0.03,0.1\) y
\(37/3\,600\) en \(0.3\)\\
\texttt{c590}, \(q=0.9999\)
& Superficie esférica
& Separación observada hasta \(r=0.3\); revisión macroscópica en \(r=1,2\)
& \(0/8\,400\) hasta \(0.3\);
\(22/2\,400\) en \(1\);
\(588/2\,700\) en \(2\), con 131 divergencias\\
\bottomrule
\end{tabularx}
\end{table}

Dentro de las bolas, horizontes y reglas de contacto del contrato local
\(r\leq0.01\), el conjunto regular no suave es compatible con la definición
local de atractor oculto. Los radios \(0.03\), \(0.1\) y \(0.3\) describen una
extensión macroscópica de la geometría de cuenca y no amplían la afirmación de
ocultedad más allá del contrato local. Para \texttt{c590}, la compatibilidad
local se limita a las superficies, radios, horizontes y reglas de contacto con
\(r\leq0.3\). Los contactos observados en \(r=1\) y \(r=2\) describen acceso
macroscópico a la cuenca. Todos estos resultados proceden de muestras finitas
y no constituyen una prueba global de separación de cuencas.
